\FloatBarrier
\section{Results}
\label{sec:results}

\subsection{Router Consolidation}
\label{sec:consolidation}

Figure~\ref{fig:consolidation} shows routing metrics as a function of $\lambda$.
All four panels respond monotonically to entropy pressure across the sweep range
$\lambda \in [0, 0.007]$, as summarized in Table~\ref{tab:consolidation}.

\begin{table}[ht]
\centering
\caption{Router consolidation metrics across the $\lambda$ sweep. All values are measured
on 40 held-out batches via \texttt{eval\_compare.py}.}
\label{tab:consolidation}
\small
\begin{tabular}{lcccc}
\toprule
$\lambda$ & $H(\mathbf{p})$ & Unique cores & Hard overlap & K=24 LRU (KB/tok) \\
\midrule
0.000 (\textsc{ctrl}) & 3.846 & 25{,}853 & 0.251 & 1{,}103 \\
0.001                 & 3.815 & 25{,}399 & 0.254 & 1{,}098 \\
0.003                 & 3.799 & 23{,}077 & 0.262 & 1{,}085 \\
0.005                 & 3.732 & 21{,}273 & 0.266 & 1{,}020 \\
0.007                 & 3.647 & 18{,}817 & 0.278 & 1{,}009 \\
\bottomrule
\end{tabular}
\end{table}

\begin{figure}[t]
\centering
\includegraphics[width=\linewidth]{fig_router_consolidation.png}
\caption{\textbf{Router consolidation metrics as a function of entropy pressure
($\lambda \times 10^3$).}
Upper left: router entropy $H(\mathbf{p})$ (lower = more concentrated pre-selection
distribution). Upper right: unique core combinations per evaluation run (lower = smaller
routing vocabulary). Lower left: K=24 LRU bytes/token (lower = better cache efficiency).
Lower right: hard-overlap Jaccard across consecutive tokens (higher = more stable routing
paths). All four panels respond monotonically. Dashed vertical lines mark the three
qualitative regimes: generalization-balanced ($\lambda{=}0.003$, blue), cache sweet spot
($\lambda{=}0.005$, orange), and boundary ($\lambda{=}0.007$, red).}
\label{fig:consolidation}
\end{figure}

Router entropy decreases monotonically (more concentrated pre-selection distribution),
unique core combinations decrease (fewer distinct routing vocabularies), and hard-overlap
Jaccard across consecutive tokens increases (more stable routing paths).
The LRU cache metric at $K = 24$ decreases from 1,103 to 1,009\kb{} across the sweep ---
a $8.5\%$ reduction in simulated bytes-in-motion.

The response is non-uniform across the range.
At $\lambda = 0.001$ the effect on routing vocabulary (unique cores) is marginal ($-1.8\%$);
the main consolidation occurs between $\lambda = 0.001$ and $\lambda = 0.005$.
At $\lambda = 0.007$ the unique-core count continues to decrease, but as
Section~\ref{sec:pareto} shows, quality effects become less favorable.
We identify three qualitative regimes: a \emph{generalization-balanced} region
($\lambda \leq 0.003$), a \emph{cache sweet spot} ($\lambda \approx 0.005$), and a
\emph{boundary} ($\lambda = 0.007$) where consolidation is maximal but the generalization
picture becomes less reliable.

\subsection{Cache--Quality Pareto}
\label{sec:pareto}

Figure~\ref{fig:pareto} plots each model as a point in the (K=24 LRU bytes/token,
code-ratio held-out) plane.
The \emph{code ratio} is defined as the ratio of held-out code loss to the mean of
held-out Web, Wikipedia, and Literature losses; lower values indicate better code
generalization relative to the other domains.

\begin{figure}[t]
\centering
\includegraphics[width=0.85\linewidth]{fig_entropy_pareto.png}
\caption{\textbf{Cache--quality Pareto under entropy pressure.}
Each point represents a model evaluated on K=24 LRU bytes/token (x-axis; lower = better
cache locality, axis reversed) and held-out code-ratio (y-axis; lower = better code
generalization).
All points share a single \textsc{ctrl} anchor from the $\lambda{=}0.003$ evaluation run
for visual consistency.
The shaded band highlights the dominance zone of $\lambda{=}0.003$ over \textsc{ctrl}.
Target-entropy variants (H=3.75, H=3.70; hollow markers) fall below the Pareto frontier
of the unconstrained sweep.}
\label{fig:pareto}
\end{figure}

\paragraph{$\lambda{=}0.003$ Pareto-dominates \textsc{ctrl}.}
In seed~0, $\lambda = 0.003$ improves K=24 cache efficiency by $1.6\%$ ($1{,}103 \to
1{,}085\kb$) while simultaneously improving the code-generalization ratio from $0.886$ to
$0.866$ ($\Delta = -0.021$).
This improvement replicates in seed~1, where the code-ratio improvement is larger
($0.896 \to 0.822$, $\Delta = -0.074$).
$\lambda{=}0.003$ is therefore a Pareto improvement over the unregularized baseline: mild
entropy pressure sharpens routing in a way that is simultaneously beneficial for cache
locality and code generalization.

\paragraph{$\lambda{=}0.005$: stronger cache, seed-dependent quality.}
$\lambda = 0.005$ achieves the largest consistent cache improvement in the sweep
($K{=}24\colon 1{,}103 \to 1{,}020\kb$, $-7.5\%$; $K{=}16\colon -5.2\%$).
Within-run paired comparisons show opposing signals across seeds: in seed~0, the code-ratio
is $0.877$ versus the within-run \textsc{ctrl} of $0.859$ ($\Delta = +0.018$, slight
degradation); in seed~1, the ratio is $0.841$ versus within-run \textsc{ctrl} $0.863$
($\Delta = -0.022$, improvement).
The Pareto figure uses a single \textsc{ctrl} anchor ($0.886$) for visual consistency;
$\lambda{=}0.005$ ($0.877$) and $\lambda{=}0.003$ ($0.866$) both fall below that anchor.
The code-quality effect of $\lambda = 0.005$ is seed-dependent and should not be
characterized as a consistent cost or benefit.

\paragraph{$\lambda{=}0.007$: cache-best boundary.}
At $\lambda = 0.007$ the cache metric reaches its lowest value in the sweep
($K{=}24\colon 1{,}009\kb$, ${\approx}1.1\%$ improvement over $\lambda{=}0.005$), while
the code-ratio returns toward \textsc{ctrl} levels ($0.884$ in seed~0).
$\lambda{=}0.007$ is cache-best but the additional gain over $\lambda{=}0.005$ is small
relative to the loss of the balanced generalization behavior observed at $\lambda{=}0.003$.
We treat $\lambda{=}0.007$ as the boundary point of the controllable axis: further entropy
pressure yields diminishing cache returns with less predictable quality effects.

\paragraph{Target-entropy variants.}
We also evaluate a bounded entropy variant,
$\lambda \cdot \mathrm{relu}(H(\mathbf{p}) - H_{\mathrm{target}})$, with targets
$H = 3.75$ and $H = 3.70$.
Both fall below the Pareto frontier of the unconstrained sweep: the bounded objective
converges more slowly within the 2,000-step continuation window and does not reach the
same cache/quality operating points.

\subsection{Dense Quality Baseline}
\label{sec:dense}

Figure~\ref{fig:quality} shows held-out loss across all four domains for the dense
baseline and three sparse models, evaluated on the same batches.
Dense d24 is the quality upper bound in every domain (Table~\ref{tab:quality}).

\begin{table}[t]
\centering
\caption{Held-out language modeling loss per domain (paired evaluation; same batches per
run, dense as consistent anchor). Dense d24 is the quality upper bound in all domains.}
\label{tab:quality}
\small
\begin{tabular}{lcccc}
\toprule
Model & Web & Wiki & Code & Lit \\
\midrule
Dense d24                        & 5.448 & 5.456 & 4.621 & 4.735 \\
SR-Core \textsc{ctrl}            & 5.807 & 5.841 & 4.856 & 5.286 \\
$+$ entmin $\lambda{=}0.003$     & 5.672 & 5.688 & 4.969 & 5.352 \\
$+$ entmin $\lambda{=}0.005$     & 5.887 & 5.792 & 4.985 & 5.419 \\
\bottomrule
\end{tabular}
\end{table}

\begin{figure}[t]
\centering
\includegraphics[width=\linewidth]{fig_dense_vs_sparse_quality.png}
\caption{\textbf{Dense quality baseline vs.\ SR-Core sparse variants.}
Held-out language modeling loss across four domains for Dense d24 and three SR-Core sparse
models, evaluated on paired held-out batches (dense as consistent anchor per run).
Dense d24 is the quality upper bound in every domain.
The $\Delta$ annotation on the Code bars shows the \textsc{ctrl}--dense gap ($\Delta = 0.24$
nats); Code has the smallest dense--sparse gap and Literature the largest ($0.55$ nats).
Entropy-minimized variants are evaluated for systems behavior (cache, routing), not
dense-equivalent quality.}
\label{fig:quality}
\end{figure}

The dense advantage is largest in Literature ($0.55$ nats) and smallest in Code
($0.24$ nats).
Entropy-minimized variants do not narrow the dense--sparse quality gap; the objective
shapes routing structure without a commensurate language modeling quality gain in paired
evaluation.
The contribution of entropy minimization is a systems-level degree of freedom, not a path
toward dense-equivalent quality.

These results are evaluated with the dense baseline as a consistent anchor (same held-out
batches for all models in each paired run).
Within-run deltas are reliable; absolute cross-run comparisons should be interpreted with
the batch-sampling variance of approximately $\pm 0.05$ nats in mind.

\subsection{Negative Controls}
\label{sec:negcontrols}

The \emph{softfull} objective --- a soft-Jaccard penalty encouraging pairwise router
similarity across consecutive tokens --- moves routing metrics in the opposite direction
from entropy minimization.
Compared to \textsc{ctrl}, \emph{softfull} increases unique core combinations from
$25{,}684$ to $30{,}246$ ($+17.8\%$) and decreases hard-overlap Jaccard from $0.259$ to
$0.251$.
Router entropy rises slightly ($3.843 \to 3.868$), and K=24 LRU bytes change by less than
$1.4\%$.

This confirms the mechanism described in Section~\ref{sec:intervention}: cross-token
similarity objectives act on pairwise overlap but do not control the within-token
distribution concentration.
The \emph{softfull} result is not a failure of regularization strength --- it is a
directional failure.
Maximizing cross-token similarity can be achieved by the router increasing route diversity
(unique routing combinations) while still mapping similar tokens to similar paths on
average.
Pre-selection entropy is the correct intervention point.

\subsection{Offloading Simulation}
\label{sec:offload}

To bridge the cache-simulation metrics to a bytes-in-motion estimate, we run the LRU
offloading simulator on the fully-trained 17k-step checkpoints (Table~\ref{tab:offload}).

\begin{table}[t]
\centering
\caption{Simulated bytes-in-motion under LRU caching at 16\,GB/s bandwidth (17k-step
checkpoints). We report $K = 8$ and $K = 16$; the dense model crossover occurs at
$K \geq 24$ (outside these columns).}
\label{tab:offload}
\small
\begin{tabular}{lccc}
\toprule
Model & K=8 KB/tok & K=16 KB/tok & vs.\ Dense (K=8) \\
\midrule
Dense d24                    & 12{,}348 & 12{,}348 & $1.00\times$ \\
SR-Core \textsc{ctrl}        & 3{,}035  & 1{,}857  & $0.246\times$ \\
$+$ entmin $\lambda{=}0.003$ & 3{,}015  & 1{,}808  & $0.244\times$ \\
$+$ entmin $\lambda{=}0.005$ & 2{,}973  & 1{,}747  & $0.241\times$ \\
\bottomrule
\end{tabular}
\end{table}

At $K = 8$ cached blocks, sparse SR-Core requires $4.1\times$ fewer bytes per token than
dense layer-offloading; entropy consolidation at $\lambda = 0.005$ adds a further $2\%$
reduction ($3{,}035 \to 2{,}973\kb$).
These are simulated estimates under an LRU caching model at 16\,GB/s bandwidth; they do not
account for transfer scheduling, compute overlap, or dispatch overhead.

One structural boundary is worth noting: at cache capacities
$K \geq n_{\mathrm{layers}} = 24$ for Dense d24, the dense model fits entirely in cache
and its bytes-in-motion approach zero, while the sparse model still incurs misses from a
bank of $n = 64$ blocks.
We report only $K = 8$ and $K = 16$ in Table~\ref{tab:offload}; the dense crossover occurs
outside these reported columns at $K \geq 24$.
This capacity crossover is not a contradiction of the sparse working-set result; it reflects
that a cache large enough to hold the entire dense model eliminates its transfer cost, an
assumption that fails precisely in the memory-constrained deployment regime this work
targets.



